\emph{This section describes the terrain in the first model generation (1,103 buildings); the version with the full city (42,018 Overture Maps footprints, channels carved according to the official hydrographic network and about 852,000 triangles) is described in Section~\ref{sec:res-v4}.} The terrain without buildings or vegetation, i.e., the digital terrain model (DTM), is FABDEM V1-2 \citep{hawker2022,fabdem2023}, with a resolution of \SI{1}{\arcsecond} (about \SI{30}{m}), bilinearly resampled to a regular \SI{4}{m} grid in the UTM system, zone 19 south (WGS~84 datum). Resampling increases point density but does not improve the effective resolution of the elevation model. This limitation must be considered when interpreting the results. A total of \num{1107} building footprints from OpenStreetMap (OSM; \citealp{osm2026}) were incorporated as raised blocks. Each block's height is the mean of $\mathrm{DSM}-\mathrm{DTM}$, where the digital surface model (DSM) is Copernicus GLO-30 \citep{copernicus_glo30} and the DTM is FABDEM, when the footprint covers at least \SI{900}{m^2} and the difference exceeds \SI{2.5}{m}. Otherwise, the height is the OSM number of stories multiplied by \SI{3}{m}; without that information, \SI{6}{m} is assumed. The mean and 95th percentile of $\mathrm{DSM}-\mathrm{DTM}$ are reported for each block. Four OSM polygons tagged as buildings that correspond to facilities (sports complexes and a shopping center with mean $\mathrm{DSM}-\mathrm{DTM}$ below \SI{4}{m}) were excluded because, represented as solid blocks in the model, they would have artificially diverted the flow. Streets were lowered by \SI{0.15}{m}, the assumed curb height, with widths based on their category.
